\section{Round-three positive closure: an exact Hamiltonian impact realization}
\label{sec:r3-a1}

This section is controlling whenever the words \emph{mechanical},
\emph{collision}, \emph{path law}, or \emph{endogenous field} are used.  It
adds an independently specified Hamiltonian impact system whose Poincar\'e map
is the four-branch family, places all parameter values on one symbolic path
space, strengthens the conditioning statement to a random-root process
principle, and states the exact axiom under which the canonical field is also
the coefficient of the excess-pressure semigroup.

\subsection{One common path space}

Put
\[
 \Sigma=\{1,2,3,4\}^{\mathbb Z},\qquad
 (\sigma\omega)_k=\omega_{k+1},
\]
and let \(\nu_a\) be the Bernoulli probability with one-symbol weights
\(w(a)\).  For a point which never meets a seam, let
\(\mathfrak c_a(x,y)\in\Sigma\) be its complete branch itinerary.

\begin{theorem}[Common-path factor and exact change of law]
\label{thm:r3-a1-common}
For every \(a\in(-1,1)\),
\[
 (\mathfrak c_a)_\#m=\nu_a .
\]
Let \(\theta\in\mathbb R\), put \(a=\Lambda'(\theta)\), and let
\(F:\Sigma\to\mathbb R\) be bounded and measurable with respect to a finite
coordinate block.  Then, for every \(n\) containing that block,
\[
 \frac{\mathbb E_{\nu_{a_0}}
  \exp\{\theta A_n+\theta F\}}
 {\mathbb E_{\nu_{a_0}}\exp\{\theta A_n\}}
 =\mathbb E_{\nu_a}e^{\theta F}.
\]
For square-coordinate observables \(G_{a_0}\) and \(G_a\), the corresponding
identity holds exactly when there is one measurable \(F\) on \(\Sigma\) with
\(G_b=F\circ\mathfrak c_b\) almost everywhere for
\(b\in\{a_0,a\}\).
\end{theorem}

\begin{proof}
The inverse image of a cylinder
\([i_{-r}\ldots i_s]\) is, modulo seams, a rectangle whose Lebesgue measure is
\(\prod_{k=-r}^s w_{i_k}(a)\).  This proves the factor identity on cylinders
and hence on the generated sigma-field.  On the common symbolic space,
\[
 \frac{d\nu_a^{[0,n-1]}}{d\nu_{a_0}^{[0,n-1]}}
 =\exp\{\theta A_n-n\Lambda(\theta)\};
\]
integration gives the second assertion.  The final assertion is precisely the
condition that the two coordinate expressions represent the same random
variable before the Radon--Nikodym derivative is applied.
\end{proof}

\subsection{Autonomous Hamiltonian impact network}

Let \((q,p)\) be canonical coordinates on the section
\(\mathcal S=\mathbb T^2\) with symplectic form \(dq\wedge dp\).  The incoming
port \(i\) is
\[
 \mathcal S_i(a)=I_i(a)\times\mathbb T.
\]
On this port define the exact affine symplectic map
\[
 \mathcal R_i^a(q,p)
 =\left(
 \frac{q-s_{i-1}(a)}{w_i(a)},
 s_{i-1}(a)+w_i(a)p
 \right).
\]
It has the type-two generating function
\[
 S_i^a(q,P)
 =\frac{(q-s_{i-1}(a))(P-s_{i-1}(a))}{w_i(a)},
\]
because \(Q=\partial_P S_i^a\) and
\(p=\partial_qS_i^a\).

To realize the reset dynamically, adjoin a clock pair \((s,E)\in
\mathbb T\times\mathbb R\), with symplectic form \(ds\wedge dE\).  Between
impact layers use the autonomous Hamiltonian \(E\), so \(\dot s=1\).  In a
thin layer before \(s=1\), apply successively the Hamiltonian flows which
translate \(q\) by \(-s_{i-1}(a)\), evolve for unit time under
\[
 H_i^a(q,p)=-(\log w_i(a))qp,
\]
and translate \(p\) by \(s_{i-1}(a)\).  The layer is selected by the incoming
port \(q\in I_i(a)\).  Equivalently, a single branchwise impulsive Hamiltonian
is obtained by composing these three compactly supported time-dependent
Hamiltonians and autonomizing with the clock pair.

\begin{theorem}[Mechanical realization of the driven family]
\label{thm:r3-a1-impact}
There is a piecewise smooth autonomous Hamiltonian impact flow
\(\Phi_a^t\) on the four-dimensional extended phase space such that:
\begin{enumerate}[label=(\roman*)]
\item the flight time between successive incoming sections is one;
\item the impact hypersurface, incoming port, branch Hamiltonian, and
      canonical impact law are the objects displayed above;
\item the first-return map on \(\mathcal S\) is exactly \(B_a\);
\item Liouville measure induces normalized Lebesgue measure on
      \(\mathcal S\), and the incoming flux through port \(i\) is
      \(w_i(a)\);
\item after coupling the channel one-form
      \(\eta=c_i\,ds\) on port \(i\), the work accumulated during the first
      \(n\) flights is exactly \(A_n\).
\end{enumerate}
Consequently the canonical tilt of the \(a_0\)-impact system by
\(\theta A_n\) is the natural Liouville path law of the mechanically distinct
member \(\Phi_{\Lambda'(\theta)}^t\) of the same impact family.
\end{theorem}

\begin{proof}
The Hamiltonian translation flows are canonical.  The middle flow solves
\[
 \dot q=-(\log w_i)q,
 \qquad
 \dot p=(\log w_i)p,
\]
so its time-one map is \((q,p)\mapsto(q/w_i,w_i p)\).  Conjugating by the two
translations gives \(\mathcal R_i^a\).  Autonomization by \((s,E)\) preserves
\(dq\wedge dp+ds\wedge dE\), and the section return time is one because
\(\dot s=1\) away from a layer whose clock duration is absorbed into its
standard impulsive parametrization.  Thus the return map is the displayed
baker map and preserves section area.  The flux of section area through
\(I_i(a)\times\mathbb T\) is its width \(w_i(a)\).  Finally
\(\int_0^1c_i\,ds=c_i\), so the channel work is \(A_n\).  The final assertion
follows from Theorem~\ref{thm:r3-a1-common} and the exact formula for the
Radon--Nikodym derivative.
\end{proof}

\begin{remark}[Hard-wall regularization]
Choose a nonnegative bump \(\rho\) of integral one and replace each impulsive
layer by \(\eta^{-1}\rho((s-s_i)/\eta)H_i^a\), with disjoint layer supports.
Its scattering map converges in every \(C^r\) norm on compact subsets of a
port to \(\mathcal R_i^a\) as \(\eta\downarrow0\), while the return map and
Liouville flux converge uniformly away from seams.  Hence the impact model is
the zero-thickness limit of smooth autonomous Hamiltonian channels, not a
probability kernel imposed after the dynamics.
\end{remark}

\subsection{Random-root process conditioning}

Let \(K_n=(n+A_n)/2\), condition on \(K_n=k_n\), and extend the conditioned
word periodically to a point of \(\{-,+\}^{\mathbb Z}\).  Let \(U_n\) be
uniform on \(\{0,\ldots,n-1\}\), independent after conditioning, and root the
word at \(U_n\).  The four subtypes are then sampled independently inside
each sign with their fixed conditional probabilities.

\begin{theorem}[Process-level Gibbs conditioning at a uniform root]
\label{thm:r3-a1-root}
If \(k_n/n\to q\in(0,1)\), the law of the uniformly rooted four-symbol path
converges weakly on \(\Sigma\) to \(\nu_{2q-1}\).  More quantitatively, for a
window of radius \(r<n/2\),
\[
 \left\|
 \mathcal L\big((\omega_{U_n+j})_{|j|\le r}\mid K_n=k_n\big)
 -\nu_{a_n}^{[-r,r]}
 \right\|_{\mathrm{TV}}
 \le \frac{(2r+1)2r}{2(n-2r)}.
\]
Furthermore the empirical rooted-path measure
\[
 \frac1n\sum_{j=0}^{n-1}
 \delta_{\sigma^j\omega}
\]
converges in probability, in the cylinder weak topology, to
\(\nu_{2q-1}\).
\end{theorem}

\begin{proof}
Given the total number of positive symbols, every ordered window is a sample
without replacement.  Sequential maximal coupling with Bernoulli sampling of
parameter \(q_n=k_n/n\) gives the displayed bound; random rooting removes the
endpoint asymmetry.  The subtype laws are identical under the conditioned and
Bernoulli laws, so the same estimate holds for four symbols.  For the
empirical path measure, apply the window estimate to each cylinder and use
its conditional variance bound \(O(n^{-1})+O(r^2/n)\).  A diagonal argument
over the countable cylinder algebra gives convergence in the cylinder weak
topology.
\end{proof}

\subsection{Mechanical canonical cocycles and the coefficient}

For a bounded common-path work \(F\), define the canonical excess work
functional
\[
 \mathfrak E_\theta(F)
 =\theta^{-1}\log\mathbb E_{\nu_a}e^{\theta F},
 \qquad a=\Lambda'(\theta).
\]
The following theorem states the exact, non-circular identification used later
in the series.  It does not assert that an unrelated preference functional
must use the mechanical field.

\begin{theorem}[Uniqueness on the canonical mechanical cocycle]
\label{thm:r3-a1-calibration}
Let \((\mathfrak V_{s,t})\) be a family on bounded common-path work satisfying:
\begin{enumerate}[label=(\alph*)]
\item cash additivity, monotonicity, normalization, and the exact tower law;
\item its likelihood-ratio cocycle is the mechanically prepared cocycle,
\[
 \frac{d\mathbb Q}{d\nu_{a_0}}\Big|_{\mathcal F_t}
 =\exp\{\theta A_t-t\Lambda(\theta)\};
\]
\item for every bounded increment \(F\), exponentiating the value gives a
      positive linear conditional functional under that same \(\mathbb Q\);
\item work and current are transformed by the same unit change.
\end{enumerate}
Then
\[
 \mathfrak V_{s,t}(F)
 =\theta^{-1}\log
 \mathbb E_{\nu_a}[e^{\theta F}\mid\mathcal F_s]
\]
for \(\theta\ne0\), with the conditional-expectation limit at \(\theta=0\).
In particular the coefficient on this mechanically specified cocycle is the
conjugate field \(I'(a)\).
\end{theorem}

\begin{proof}
By (b), the only conditional probability entering (c) is \(\nu_a\).  Cash
additivity implies that the positive linearization has the form
\(T_{s,t}(e^{\beta F})\) for one scalar \(\beta\) independent of the time
split; the tower law makes the scalar constant across concatenations.
Applying the cocycle identity to the generator \(A_t-A_s\) shows that its
Radon--Nikodym exponent is both \(\beta(A_t-A_s)\) and
\(\theta(A_t-A_s)\).  Since this increment takes both signs with positive
probability, \(\beta=\theta\).  Normalization fixes the logarithmic constant.
The unit covariance in (d) excludes a hidden rescaling of the work bundle.
\end{proof}

\subsection{Defined all-order response arrays}

Let \(Q_a=I-\Pi_a\) and
\(S_a=(I-L_a)^{-1}Q_a=\sum_{n\ge0}L_a^nQ_a\) on the centered Sobolev ladder.
For integers \(j\ge0\), \(k_1,\ldots,k_j\ge1\), and
\(n_0,\ldots,n_j\ge0\), define the response word
\[
 W_{\boldsymbol k,\boldsymbol n}(a)
 =L_a^{n_0}(\partial_a^{k_1}L_a)L_a^{n_1}\cdots
  (\partial_a^{k_j}L_a)L_a^{n_j}Q_a.
\]
For a bounded functional \(\ell\) on \(W^{r,1}\), the bilateral array is
\(\ell W_{\boldsymbol k,\boldsymbol n}(a)h\).

\begin{proposition}[Absolute summability and finite differences]
\label{prop:r3-a1-array}
If \(|\boldsymbol k|=\sum k_i\le q\), then, uniformly on
\(K\Subset(-1,1)\),
\[
 \sum_{\boldsymbol n\in\mathbb N^{j+1}}
 \|W_{\boldsymbol k,\boldsymbol n}(a)\|_{W^{r+q,1}\to W^{r,1}}
 \le C_{K,r,q,j}(1-\rho_K)^{-(j+1)}.
\]
The same estimate holds after application of \(\ell\).  Every derivative of
\(S_a\) is a finite sum of these absolutely convergent arrays, and the order
\(q\) finite-difference remainder is bounded by
\(C|a-b|^{q+1}\) from \(W^{r+q+1,1}\) to \(W^{r,1}\).
\end{proposition}

\begin{proof}
A parameter derivative of order \(k_i\) loses at most \(k_i\) Sobolev
orders.  Between derivative insertions, the centered transfer operator has
norm at most \(\rho_K^{n_i}\).  Multiplication of the insertion bounds and
summation of the \(j+1\) independent geometric series gives the estimate.
Differentiate \((I-L_a)S_a=Q_a\) recursively to obtain finite sums of the
same words.  Taylor's formula with integral remainder and the estimate at
order \(q+1\) gives the finite-difference assertion.
\end{proof}

\begin{theorem}[Round-three A1 closure]
\label{thm:r3-a1-main}
The four-branch path ensemble is simultaneously:
\begin{enumerate}[label=(\roman*)]
\item the Liouville return process of an explicitly specified autonomous
Hamiltonian impact network;
\item an exact exponential family on one common symbolic path space;
\item the process-level limit of the uniformly rooted microcanonical current
ensemble;
\item a mechanically driven family under canonical tilting; and
\item the base of a uniquely calibrated excess-pressure cocycle when the
canonical-cocycle axioms of Theorem~\ref{thm:r3-a1-calibration} are imposed.
\end{enumerate}
All response arrays used downstream are the operators defined above and are
absolutely summable on the stated Sobolev ladder.
\end{theorem}

\begin{proof}
Combine Theorems~\ref{thm:r3-a1-common},
\ref{thm:r3-a1-impact}, \ref{thm:r3-a1-root},
\ref{thm:r3-a1-calibration}, and
Proposition~\ref{prop:r3-a1-array}.
\end{proof}
